\documentclass[twocolumn,12pt]{article}    % two-column because cheat sheets fit like that
\usepackage{preamble}
\newcommand\coursename{Intro to Optimization}
\newcommand\Def{\text{\textbf{Def.}} \hspace{5pt}}
\newcommand \rn {\mathds{R}^n}
\newcommand \hfive {\hspace{5pt}}

\begin{document}
\footnotesize % small because, it's a cheat sheet after all, why not cram it
\subsection*{1st Half}
\begin{equation*}
    \|\cdot\| : \mathds{R}^n \hspace{-3pt} \to \hspace{-3pt} \mathds{R} \, s. t:\, \|\Vec{x}\| \geq 0, \, \|\lambda \Vec{x}\| = |\lambda| \|\Vec{x}\|, \, \|\Vec{x}+\Vec{y}\| \leq \|\Vec{x}\| +\|\Vec{y}\|
\end{equation*}
\begin{equation*}
    |\Vec{x}^T\Vec{y}| \leq \|\Vec{x}\|_2 \|\Vec{y}\|_2   \, (C.S.)\hspace{7pt} \mathds{R}^n_+ = \{\mathds{R}^n \geq0\}  \hspace{7pt} \mathds{R}^n_{++} = \{\mathds{R}^n > 0\}
\end{equation*}
\begin{equation*}
    \Delta_n  = \{\Vec{x} \in \mathds{R}^n : x_i\geq 0 \forall i \in \{1,...,n\},\, \Vec{1}^T\vec{x} = 1\}
\end{equation*}
\begin{equation*}
    B(\vec{c},\,r) = \{\vec{x} \in \mathds{R}^n : \|\vec{x} -\vec{c}\| < r\} \hspace{20pt} \text{// open( ), closed[ ]}
\end{equation*}

\vspace{-20pt} 
\begin{equation*}
    \left\|\Vec{x}\right\|_p \triangleq \left( \sum^n_{i=1} |x_i|^p\right)^{1/p} \hspace{10pt}  \|A\|_F =  \sqrt{\Sigma_{i}\Sigma_{j} |A_{ij}|} 
\end{equation*}
\begin{equation*}
    \|A\Vec{x}\|_b \leq \|A\|_{a,b} \|\Vec{x}\|_a \hspace{7.5pt} \text{w/} \hspace{7.5pt} \ \|A\|_{a,b} = \max_{\Vec{x} \in \mathds{R}^n} \{\|A\Vec{x}\|_b : \|\Vec{x}\|_1 \leq 1\}
\end{equation*}
\begin{equation*}
    \|A\|_2 = \|A\|_{2,2} = \sqrt{\lambda_{\max}(A^TA)} = \sigma_{\max}(A) \hspace{5pt} \text{// spectral norm}
\end{equation*}
\begin{equation*}
    \|A\|_1 = \max_{i=1,...,n} \left\{\text{col. sums}\right\} \hspace{19pt} \|A\|_\infty = \max_{j=1,...,n} \left\{\text{row. sums}\right\}
\end{equation*}
\begin{equation*}
      R_{A}(\vec{x}\neq 0 ) =\frac{\vec{x}^TA\vec{x}}{\vec{x}^T\vec{x}}\hspace{30pt} \lambda_{\min}(A) \leq R_{A}(\vec{x}\neq0)\leq \lambda_{\max}(A) 
\end{equation*}
\begin{equation*}
      \lambda_{\min}(A) = \min_{\vec{x}\neq0} R_{A}(\vec{x}) \hspace{19pt} \lambda_{\max}(A) = \max_{\vec{x}\neq0} R_{A}(\vec{x})
\end{equation*}
\begin{equation*}
     f'(\vec{x};\, \vec{d}_i) = \lim_{t\to0^+} \frac{f(\vec{x} +t \vec{d}_i) -f(\vec{x})}{t} = \nabla f(\vec{x})^T \vec{d} \hspace{5.5pt} \text{// dir. deriv.}
\end{equation*}
\begin{equation*}
    \nabla f(\Vec{x}) = \begin{bmatrix} \frac{\partial{f}}{\partial{x_1}}\\ \vdots \\ \frac{\partial{f}}{\partial{x_n}} \end{bmatrix} 
    \in \mathds{R}^n \hspace{10pt}
    \nabla^2 f(\vec{x}) = \begin{bmatrix}
        | & & |\\
        \nabla f_1 & \dots & \nabla f_n\\
        | & & |\\
    \end{bmatrix} \in \mathds{R}^{n \times n}
\end{equation*}
\begin{align*}
    \text{TS} \hfive f(\vec{y}) = &f(\vec{x}) + \nabla f(\vec{x})^T(\vec{y}-\vec{x}) \\
    &+ \underbrace{\frac{1}{2} (\vec{y}-\vec{x})^T \nabla^2 f(\Vec{x}) (\vec{y}-\vec{x})}_{=o\|\vec{y}-\vec{x}\| \text{ for lin. approx}} + \underbrace{o\left( \|\vec{y} -\vec{x}\| ^2 \right)}_{\text{quad. approx}}
\end{align*}
\begin{equation*}
    f(\vec{x}^* + \vec{d}) = f(\vec{x}^*) + \hspace{-5pt}
    \underbrace{\nabla f(\vec{x}^*)\vec{d}}_{= 0 \text{ if stat. pt.}} \hspace{-5pt} + \frac{1}{2}\vec{d}^T\nabla^2f(\vec{x}^*) \vec{d} + o (\|\vec{d}\|^2 )
\end{equation*}
\begin{equation*}
    \lim_{\|\vec{x}\| \to \infty} f(\vec{x}) =\infty \implies \nabla^2f(\vec{x}) \succ 0 \hspace{25pt} \text{// Coercive} 
\end{equation*}
\noindent \underline{Necessary 2\(^{\text{nd}}\) Order Optimality}
\begin{equation*}
    \text{if } \vec{x}^* \text{ is \textbf{local} min, then } \nabla^2f(\vec{x}^*) \succeq 0 \hspace{10pt} \textbf{local} \max \text{if} \preceq 0
\end{equation*}
\noindent \underline{Sufficient 2\(^{\text{nd}}\) Order Optimality}
\begin{equation*}
    \text{if } \nabla^2 f(\vec{x}^*) \succ 0 \text{ then } \vec{x}^* \text{ is \textbf{strict local} min} \hspace{10pt} \max \text{if} \prec 0
\end{equation*}
\begin{equation*}
    \text{if } \nabla^2 f(\vec{x}^*) \text{ is indef. then } \vec{x}^* \text{ is a \textbf{saddle}}
\end{equation*}
\noindent\underline{Global Optimality}
\begin{equation*}
    \text{if } \nabla^2f(\vec{x}^*) \succeq 0,\, \forall \vec{x},\,\vec{x}^*= \text{stat. pt, then } \vec{x}^* \text{ is \textbf{global} min} 
\end{equation*}
\begin{equation*}
    \text{if } \exists m>0 \, \text{s.t. } \nabla^2 f(\vec{x}) \succ mI, \, \forall \vec{x} \text{ then:} 
\end{equation*}
\begin{equation*}
    f(\vec{x}) \geq f(\vec{x}^*) + \frac{m}{2} \underbrace{ \|\vec{x} - \vec{x}^*\|^2}_{\text{bowl above } f(\vec{x}^*)} \geq f(\vec{x}^*)
\end{equation*}\vspace{-10pt}
\begin{equation*}
    \overbrace{A\vec{x} = \vec{b} \hspace{20pt} A\in\mathds{R}^{m\times n } }^{\text{Least Squares}} \hspace{20pt} \begin{cases}
        &m>n\, \text{overdetermined } \\
        &m<n \, \text{underdetermined} 
    \end{cases}
\end{equation*}\vspace{-10pt}
\begin{equation*}
    \min_{\vec{x} \in \mathds{R}^n} \overbrace{\|A\vec{x} - \vec{b}\|_2^2}^{\text{residuals}} \implies \vec{x}^TA^T A\vec{x} - 2\vec{b}^TA\vec{x} - \vec{b}^T\vec{b}
\end{equation*}
\begin{equation*}
    \nabla_x(\cdot) =0 \implies A^TA \vec{x} = A^T\vec{b} \hspace{30pt} \text{// Normal eqs}
\end{equation*}
\begin{equation*}
    \text{if } m>n,\, A^TA \succ 0 \implies \vec{x} = (A^TA)^{-1} A^T \vec{b}
\end{equation*}
\begin{equation*}
    f(x) = ax^2 +bx +c\implies \begin{bmatrix}
        x_1^2 & x_1 & 1 \\
        \vdots & \vdots & \vdots \\
        x_m^2 & x_m & 1
    \end{bmatrix} \begin{bmatrix}
        a\\b\\c
    \end{bmatrix} = \begin{bmatrix}
        y_1 \\ \vdots \\ y_m
    \end{bmatrix}
\end{equation*}
\noindent\underline{Gradient Method} \hspace{100pt} Converges linearly
\begin{equation*}
    \textbf{def. } \vec{d} \in \mathds{R}^n \text{ is descent dir. of } f \text{ at } \vec{x} \text{ if } \begin{cases}
        &\nabla f(\vec{x})^T\vec{d} <0\\
        & \nabla f(\vec{x};\,\vec{d}) <0
    \end{cases}
\end{equation*}
\begin{equation*}
    \begin{matrix}
        \text{gradient w/ exact}\\
        \text{line search satisfies}
    \end{matrix} \, f(\vec{x}_{k+1}) \leq \left(\frac{\kappa -1}{\kappa+1}\right)^2 f(\vec{x}_k),\hspace{5pt}
    \kappa = \frac{\lambda_{\max}(A)}{\lambda_{\min}(A)}
\end{equation*}
\begin{equation*}
    \text{if } f\in C^1(\mathds{R}^n), \text{ bounded below, and} \nabla f(\vec{x}) \in C_L^{1,1},\, L,M>0:
\end{equation*}
\begin{equation*}\begin{matrix}
    f \text{ flattens}\\ \text{ as }k\to\infty \hspace{10pt}
\end{matrix}
    \begin{cases}
    & f(\vec{x}_k) - f(\vec{x}_{k+1}) \geq M\| \nabla f(\vec{x}_k) \|^2 \geq 0 \\
    & \min_{k=0,..n} \| \nabla f(\vec{x}_k) \| \leq \sqrt{\frac{f(\vec{x}_0) - f^*}{M(n+1)}}
\end{cases}
\end{equation*}
\noindent\underline{Descent Direction "Method"}\\
Inputs: \(\vec{x}_0 \in \mathds{R}^n\), tolerance \(\varepsilon >0\)\\
While \(\|\nabla f(\vec{x})\| > \varepsilon\)
\begin{enumerate}
    \item Choose \(\vec{d}_k\)
    \begin{equation*}
    \hspace{-10pt} \vec{d}_k =\begin{cases}
        & - \nabla f(\vec{x}) \hspace{42.5pt} \text{// gradient/steepest descent}\\
        & -[\nabla^2 f(\vec{x})]^{-1} \nabla f(\vec{x}) \,\text{// Newton (only if }H\succ 0)
    \end{cases}
    \end{equation*}

    \item Choose step size \(t_k\) via line search
    \begin{equation*}
        t_k = \begin{cases}
            &\tilde{t} \hspace{102.5pt} \text{// constant}\\
            &\arg\, \min_{t\geq0} f(\vec{x}_k + t_k\vec{d}_k) 
            \hspace{20pt} \text{// exact}\\
            & \hspace{87.5pt}\text{// see backtrack}
        \end{cases}
    \end{equation*}

    \item Set \(\vec{x}_{k+1} = \vec{x}_k + t_k \vec{d}_k\), \(k=k+1\)
\end{enumerate}
\noindent \underline{Backtracking Line Search}\\
Inputs: \(s>0,\, \alpha \in (0,\,1),\, \beta\in (0,\,1)\)
\begin{enumerate}
    \item Set \(t_k =s\)
    \item While \(f(\vec{x}_k +t_k\vec{d}_k) > f(\vec{x}_k) + \alpha t_k \nabla f(\vec{x}_k)^T \vec{d}_k\)\\
    \hspace{80pt} \(t_k = \beta t_k\)
\end{enumerate}
    
\noindent\underline{Lipschitz Continuity}\\
"If the next derivative is bounded by some L for all x, then the function is Lipschitz at the current derivative with L being the boundary".
\begin{equation*}
    \textbf{def. } f \in C_L^{1,1} (\mathds{R}^n)\, \Longleftrightarrow \,\|\nabla f(\vec{x}) - \nabla f(\vec{y}) \| \leq L\|\vec{x} -\vec{y}\| \hspace{20pt} 
\end{equation*}
\begin{equation*}
    \text{thm }4.20 \hspace{10pt} f \in C_L^{1,1} (\mathds{R}^n)\, \Longleftrightarrow\, \|\nabla^2 f(\vec{x})\| \leq L, \, \forall \vec{x} \in \mathds{R}^n
\end{equation*}
\begin{equation*}
    C^n < C^{a,b}_L < C^{n+1} \hspace{30pt} \text{// Lipschitz is  "in between"}
\end{equation*}
\begin{equation*}
    f \in C_L^{1,1} (\mathds{R}^n) \implies f(\vec{y}) \leq f(\vec{x}) + \nabla f(\vec{x})^T (\vec{y} -\vec{x}) + \frac{L}{2} \|\vec{y} - \vec{x}\|^2
\end{equation*}
\begin{equation*}
    f \in C_L^{1,1} (\mathds{R}^n) \implies f(\vec{x}) - f(\vec{x} -t\nabla f(\vec{x})) \geq t(1-Lt/2)\|\nabla f(\vec{x})\|^2
\end{equation*}
\noindent \underline{(Pure) Newton's Method} \hspace{40pt} Converges quadratically\\
Inputs: \(\vec{x}_0\in \mathds{R}^n,\) tolerance \(\varepsilon >0\)\\
While \(\|\nabla f(\vec{x})\| > \varepsilon\)
\begin{enumerate}
    \item Set \(\vec{d}_k= - [\nabla^2 f(\vec{x}_k)]^{-1} \nabla f (\vec{x}_k)\)
    \item Set \(\vec{x}_{k+1} = \vec{x}_k + \vec{d}_k\) \hspace{30pt} // \(t_k =1\)
    \item Set \(k = k+1\)
\end{enumerate}
"Newton's method approximates a function as quadratic and solves it directly. If the function is quadratic already, it solves it in one step".\\
"Only descends if \(A\succ0\), and still can fail. Expensive, solves linear system (\(\nabla^2\)) at every \(k\)".
\begin{equation*}
    \text{Assume } \nabla^2 f(\vec{x}) \succ mI,\, \forall\vec{x} \in \mathds{R}^n,\,\text{ then:} 
\end{equation*}
\begin{equation*}
    \begin{cases}
        &\|\nabla^2 f(\vec{x}) - \nabla^2 f(\vec{y})\| \leq L \| \vec{x} - \vec{y}\| \hspace{45pt} \text{// Lipschitz}\\
        & \|\vec{x}_{k+1} - \vec{x}^* \| \leq \frac{L}{2m} \|\vec{x}_k -\vec{x}^* \|^2 \hspace{30pt} \text{// Converge quad.}
    \end{cases}
\end{equation*}
\noindent\underline{Damped Newton's Method}\\
Newton's Method but Backtracking for step size.\\
\begin{equation*}
    \Def \text{Weierstrass Thm} : \text{if} \hfive f\in C^1(C\subseteq \rn),\, C \hfive \text{non-empty,}
\end{equation*}
\begin{equation*}
    \text{compact}: \exists \hfive \text{global (min, max) of }f \text{ over } C
\end{equation*}


%%%%%%
\subsection*{2nd Half}
\begin{equation*}
    \min \text{concave f}(\text{over convex set}) \implies \exists \min \text{ at ext pt}
\end{equation*}
In strictly convex opt prob, local min are global min \\
The set of opt pts to convex opt prob is also convex
\begin{equation*}
    \text{\textbf{Def.}} \hspace{10pt} C \text{  convex} \hspace{5pt} \iff \hspace{5pt} \vec{x},\, \vec{y},\, [\lambda \vec{x} + (1-\lambda)\vec{y}] \in C,\, \lambda \in [0,\,1]
\end{equation*}
\begin{equation*}
        \underbrace{\text{if } S \subseteq \mathds{R}^n, \, \vec{x} \in \text{conv}(S)}_{\text{Caratheodry thm}}; \hspace{10pt} \begin{matrix}
        \exists \, \vec{x}_1,\, \dots\, \vec{x}_{n+1} \in S \\
        \text{s.t.} \vec{x} \in \text{conv}(\{\vec{x}_1,\, \dots\, \vec{x}_{\textbf{n+1}}\})
    \end{matrix}
\end{equation*}
\begin{equation*}
    \text{\textbf{Def.}} \hspace{5pt} S\hspace{5pt}\text{is convex cone} \hspace{5pt} \iff \begin{cases}
        \vec{x},\, \vec{y} \in S \implies \vec{x} + \vec{y} \in S\\
        \vec{x} \in S,\, \lambda \geq 0 \implies \lambda\vec{x} \in S
    \end{cases}
\end{equation*}
\begin{equation*}
    \text{\textbf{Def.} conic comb.}\hspace{5pt} \lambda_1\vec{x}_1  + \dots + \lambda_k\vec{x}_k \hspace{5pt} \text{w/} \hspace{5pt} \begin{matrix}  \vec{x}_1,\, \dots \, \vec{x}_n \in \mathds{R}^n,\\ \vec{\lambda} \in \mathds{R}_+^k \end{matrix}
\end{equation*}
\begin{equation*}
    \Def \begin{matrix} \vec{x}\in S \hspace{5pt} \text{is \textbf{ext pt} of } \\  S \subseteq \mathds{R}^n, \text{convex}\end{matrix}\hspace{5pt} \text{if}: \hspace{5pt} \begin{matrix} \nexists \vec{x}_1 \neq \vec{x}_2 \in S,\\ \text{s.t. }  \vec{x} = \lambda\vec{x}_1 + (1-\lambda)\vec{x}_2\end{matrix}
\end{equation*}
\begin{equation*}
    \text{if} \hspace{5pt} S \subseteq \rn \hspace{5pt} \text{is compact, convex}: \hspace{5pt} S = \text{conv(ext(}S))
\end{equation*}
\begin{equation*}
    \hspace{-10pt} \Def f:C\rightarrow\rn \text{convex if}:
\end{equation*}
\begin{equation*}
    f(\lambda \vec{x} +(1-\lambda)\vec{y}) \leq \lambda f(\vec{x}) + (1-\lambda)f(\vec{y}) \hspace{20pt} \forall \vec{x},\, \vec{y} \in C,\, \lambda \in [0,\,1] 
\end{equation*}
\begin{equation*}
    \Def \text{Jenson's ineq}: \hfive f\left(\Sigma_i \lambda_i \vec{x}_i\right) \leq \Sigma_i\lambda_if(\vec{x}_i)
\end{equation*}
\begin{equation*}
    \hspace{-10pt}\begin{cases}
        f \hfive \text{convex} \iff f(\vec{x}) + \nabla f(\vec{x})^T (\vec{y}-\vec{x}) \leq f(\vec{y}) \hfive \forall \vec{x},\, \vec{y} \in C  & 1^{\text{st}}\\
        f \hfive \text{convex} \iff \nabla^2f(\vec{x}) \succeq 0 \forall \vec{x} \in C \hspace{20pt} (\text{strictly if} \succ) & 2^{\text{nd}}
    \end{cases}
\end{equation*}
\begin{equation*}
    \Def \text{Orthog. Proj:} \hfive P_c (\vec{x})  = \arg \min_{\vec{y}\in C} ||\vec{x} - \vec{y}||^2 \hspace{20pt} // \text{unique}
\end{equation*}
\begin{equation*}
    \Def \hfive \vec{x}^* \in C \hfive \text{is \textbf{stat pt} if}: \hfive \nabla f(\vec{x}^*)^T(\vec{x}-\vec{x}^*) \geq 0 \hfive \forall \vec{x}\in C 
\end{equation*}
\begin{equation*}
    \begin{cases}
        f\in C^1 \implies \hfive \text{local min are stat pts}  & \text{nec for opt}\\
        f \hfive \text{also convex}  & \text{suff for opt}
    \end{cases}
\end{equation*}
\noindent\underline{Gradient Projection Method} \hspace{30pt} Converges sublinearly\\
Set \(\vec{x}_0 \in C \), \(\varepsilon>0\), and repeat for \(k= 0,\,1,\, \dots\)
\begin{enumerate}
    \item Pick step size \(t_k\) (small, or line search)
    \item \(\vec{x}_{k+1} = P_c(\vec{x}_k - t_k\nabla f(\vec{x}_k) )\)
    \item If \(||\vec{x}_{k+1} -\vec{x}_k || \leq \varepsilon\), STOP
\end{enumerate}
Unlikely for \(\nabla f(\vec{x}^*) \approx 0\). Different from other methods. Sub-linear converge if Lipschitz \(\tilde{t}\) chosen well. (With \(n\) = \# iters)

\noindent\underline{KKT}
\begin{equation*}
    f(\vec{x}_k) - f^* \leq ||\vec{x}_0 -\vec{x}^*||^2 /(2\tilde{t} n)
\end{equation*}
\begin{align*}
    (P)\hspace{20pt} &\min_{\vec{x} \in \rn} &\,f(\vec{x}) \hspace{30pt} & \, f,\, g_i,\, h_j  \in C^1(\rn)\\
    & \hspace{10pt} \text{s.t.} &\, g_i(\vec{x}) \leq 0 \hspace{10pt} &  i = 1,\, \dots\, m\\
    & &\, h_j(\vec{x}) = 0 \hspace{10pt} & j = 1,\, \dots \, p
\end{align*}
\begin{equation*}
    \Def I(\vec{x}^*) = \{i\in [m] : g_i(\vec{x}^*) =0  \} = \hfive \text{set of active constr.}
\end{equation*}
\begin{equation*}
    \Def \text{feasible pt } \vec{x}^* \text{ is \textbf{regular} if}: 
\end{equation*}
\begin{equation*}
    \{\nabla g_i(\vec{x}^*): i \in I(\vec{x}^*)\} \bigcup \{\nabla h_j(\vec{x}) : j = 1,\, \dots\, p \} \hspace{2.5pt} \text{are lin. indep}
\end{equation*}
\begin{equation*}
    \Def \text{Slater's Condition holds if}: \hfive \exists \hat{\vec{x}} \in \rn \hfive \text{s.t.} \hfive g_i(\vec{x}) < 0 \, \forall i
\end{equation*}
\begin{equation*}
    \begin{cases}
        \text{if} \hfive \begin{matrix}
            \vec{x}^* \hfive \text{is local min of } (P)\text{, regular} \\
            \text{(Slater's cond holds)}
        \end{matrix}\implies \text{KKT pt} & \text{nec}\\
        \text{and if} \hfive (P) \hfive \text{convex} \implies \hfive \text{KKT pt IS opt pt} & \text{suf}
    \end{cases}
\end{equation*}
\noindent \underline{KKT Conditions}
\begin{equation*}
    L(\vec{x},\, \vec{\lambda},\, \vec{\mu}) = f(\vec{x})  + \Sigma \lambda_i g_i(\vec{x}) + \Sigma \mu_j h_j(\vec{x})
\end{equation*}
\begin{equation*}
    \nabla_{\vec{x}}L = \vec{0} \hspace{30pt} \lambda_ig_i(\vec{x}) = 0 \hspace{30pt} \mu_j h_j(\vec{x}) = 0
\end{equation*}
\begin{equation*}
    g_i(\vec{x}) \leq 0 \hspace{30pt} h_j(\vec{x}) = 0 \hspace{30pt}
    \vec{\lambda} \geq 0
\end{equation*}

\noindent \underline{Duality}
\begin{equation*}
    q(\vec{\lambda},\, \vec{\mu}) = \min_{\vec{x}} L(\vec{x},\, \vec{\lambda},\, \vec{\mu}) \hspace{20pt} //\text{always concave, sol is ext pt}
\end{equation*}
\begin{equation*}
    (D) \hfive q^* = \max_{\vec{\lambda},\, \vec{\mu}} q(\vec{\lambda},\, \vec{\mu}) \hspace{20pt} // \text{always convex, even if } (P) \text{ isn't}
\end{equation*}
\begin{equation*}
    \text{dom} (q) = \{ (\vec{\lambda},\, \vec{\mu}) \in \mathds{R}_+^m \times \mathds{R}^p : q(\vec{\lambda},\, \vec{\mu}) > -\infty \}
\end{equation*}
\textbf{Weak Duality}: \(q^* \leq f^*\) always holds\\
\textbf{Strong Duality}: \(q^* =f^*\) when \((P)\) convex, w/ \(\min > - \infty\)


\noindent\underline{Advanced Variants}\\
"Dream: by approximating the Hessian instead of computing, the number of iterations can be reduced w/o incurring the cost of computing the Hessian".
\begin{itemize}
    \item Diagonal scaling: Set \(\vec{d}_k = D_k \nabla f(\vec{x}_k) = SS^T\nabla f(\vec{x}_k) \) where \(\vec{x} = S\vec{y}\) is a change in variables to  alter \(x_i\)'s, or non-dimensionalize to have similar sensitivity to changes across all \(x_i\)'s.
    \item  BFGS Quasi-Newton: Alters the problem to utilize the inverse-Hessian while avoiding the issue of ever inverting. Can converge super-linearly (faster than grad. desc, slower than Newton's). Limited memory variant exist. Most opt solvers use this method.
    \item Gauss-Newton Method: Solves the non-linear least squares problem when \(f: \mathds{R}^n \to \mathds{R}^m\) via Jacobian.
\end{itemize}








%%%%
\newpage
\section*{Content that didn't fit}
\begin{equation*}
    \frac{(\vec{x}^T \vec{x})^2}{(\vec{x}^T A^{-1} \vec{x})} \geq \frac{4\kappa}{(\kappa+1)^2} \hspace{10pt}
    \kappa = \frac{\lambda_{\max}(A)}{\lambda_{\min}(A)} \hspace{10pt} \text{// Kantrovic ineq.}  
\end{equation*}
\noindent \underline{Symmetric, Square \(A\) Properties}
\begin{equation*}
     \text{Tr}(A) = \sum_{i=1}^n A_{ii} = \sum_{i=1}^n \lambda_i(A) \hspace{30pt} \text{det}(A) = \prod_{i=1}^n \lambda_i(A) 
\end{equation*}
\begin{equation*}
      R_{A}(\vec{x}\neq 0 ) =\frac{\vec{x}^TA\vec{x}}{\vec{x}^T\vec{x}}\hspace{30pt} \lambda_{\min}(A) \leq R_{A}(\vec{x}\neq0)\leq \lambda_{\max}(A) 
\end{equation*}
\begin{equation*}
    f\in C^2(\mathds{R}^n) \hspace{30pt} \left(\nabla^2 f(\vec{x}) \right)_{ij} =  \frac{\partial^2f}{\partial x_i \partial x_j} \in \mathds{R}^{n \times n}
\end{equation*}
\begin{equation*}
    h(\vec{x}) = f(g(\vec{x})) \hspace{18pt} \nabla h(\vec{x}) = \nabla g(\vec{x}) \nabla f(g(\vec{x})) \hspace{18pt} \text{// chain rule}
\end{equation*}
\noindent \underline{Sets} 
\begin{equation*}
    \text{int}(U) = \{\vec{x} \in U : B(\vec{x},\,r) \subseteq U : \text{ for some } r >0 \}
\end{equation*} 
\begin{equation*}
    \text{Cl}(U) = \bigcap \{T: T \supseteq U,\, T \text{ closed}\}
\end{equation*}
\begin{equation*}
    U\subseteq\mathds{R}^n \text{ bounded if } \exists M>0 \text{  s.t.  } U \subset B(\vec{0},\, M)
\end{equation*}
\begin{equation*}
    \text{compact} = \text{closed} + \text{bounded}
\end{equation*}
\begin{equation*}
    S \subseteq \mathds{R}^n \hspace{20pt} f^* =  \inf \{ f(\vec{x}) : \vec{x} \in S\} = \underbrace{\min \{ f(\vec{x}) : \vec{x} \in S\}}_{\text{if proper min. in set}}
\end{equation*}
\noindent \underline{Definiteness of \textbf{Symmetric} Matrices}
\begin{table}[H]
    \centering
    \footnotesize
    \begin{tabular}{ccc}
        Name & Symbolic & Eigenvalues \(\forall i\)\\
        \hline
        PSD & \(A\succeq 0\) & \(\lambda_i(A)\geq0\) \\
        PD &  \(A\succ 0\) & \(\lambda_i(A)>0\) \\
        NSD &  \(A\preceq 0\) & \(\lambda_i(A)\leq0\) \\
        ND &  \(A\prec 0\) & \(\lambda_i(A)<0\) \\
        Ind. & \(A \npreceq 0,\, A \nsucceq 0\) & \(\lambda_{\max}(A) > 0, \, \lambda_{\min}(A)<0\)\\
    \end{tabular}
\end{table}
\vspace{-30pt}%%%%%%% 
\begin{equation*}
    \textbf{Def. }A\in \mathds{R}^{n\times n} \succeq 0 \hspace{10pt} \text{if} \hspace{10pt} \vec{x}^TA\vec{x} \geq 0 \hspace{10pt} \forall \vec{x} \in \mathds{R}^n \hspace{45pt} 
\end{equation*}
\begin{equation*}
    \textbf{Def. }A\in \mathds{R}^{n\times n} \succ 0 \hspace{10pt} \text{if} \hspace{10pt} \vec{x}^TA\vec{x} \geq 0 \hspace{10pt} \forall \vec{x} \in \mathds{R}^n \setminus \{0\} \hspace{20pt}
\end{equation*}
\begin{equation*}
    A \succ 0 \rightarrow A_{ii} > 0 \hspace{30pt}  A \succeq 0 \rightarrow A_{ii} \geq 0
\end{equation*}
\begin{equation*}
    \text{if signs on diagonal differ, }A \text{ is Indef.}
\end{equation*}



\noindent\underline{Hybrid Gradient Newton Method}
\begin{equation*}
    \text{if } \nabla^2 f(\vec{x_k}) \succ 0\implies \vec{d}_k = - [\nabla^2 f(\vec{x}_k)]^{-1} \nabla f( \vec{x}_k)
\end{equation*}
\begin{equation*}
    \text{otherwise: } \vec{d}_k = -\nabla f(\vec{x}_k)
\end{equation*}
\noindent \underline{Circle Fitting}\\
"Non-linear least squares problem that can be reformulated as a linear least squares problem".
\begin{equation*}
    \text{fit most pts. } \vec{a}_i \text{ with } C(\vec{x},\, r) = \{\vec{a} \in\mathds{R}^n:\, \|\vec{a} - \vec{x}\| = r\}
\end{equation*}
\begin{equation*}
    \text{Orig prob:     } \min_{\vec{x},\, R}  \sum_{i=1}^m\left(-2\vec{a}_i^T\vec{x} + R + \|\vec{a}_i\|^2 \right)^2
\end{equation*}
\begin{equation*}
    \text{Equivalent to:    }\min_{\vec{y} \in\mathds{R}^{n+1}} \left\|\tilde{A} \vec{y} - \vec{b} \right\|^2 \hspace{5pt} \text{where:} \hspace{5pt} \vec{y} = \begin{bmatrix}
        \vec{x} \\R 
    \end{bmatrix} \in \mathds{R}^{n+1}
\end{equation*}
\begin{equation*}
    \tilde{A} = \begin{bmatrix}
        \--  2\vec{a}_1^T  \-- & -1 \\
          \vdots  &\\
         \--  2\vec{a}_m^T  \-- & -1 \\
    \end{bmatrix} \in \mathds{R}^{m\times n+1} \hspace{10pt} \Vec{b} = \begin{bmatrix}
        \|\vec{a}_1\|^2 \\ \vdots \\\|\vec{a}_m\|^2
    \end{bmatrix} \in \mathds{R}^m
\end{equation*}
\begin{equation*}
    \text{gives:} \hspace{20pt} \vec{x}^* = y^*[1:n], \hspace{5pt} r^* = \sqrt{\|\vec{x}^*\|^2 - y^*[n+1]}
\end{equation*}
\end{document}